\section{Runner reductions and projective cones}\label{sec:runners}

We now reduce the unbounded ranks in the odd-prime classical problem.
The operations will be split Harish--Chandra induction and restriction,
followed by projections onto unions of blocks.  Their matrices on the
basic rows are particularly simple.  It is these matrices, together
with positivity on projectives, that will transfer bounds.

Runner rearrangements connect this argument with the abacus
reductions of Scopes and Jost and the classical unipotent
equivalences of Hiss--Kessar
\cite{Scopes,Jost,HissKessarScopes,HissKessarScopesII}.
Here the transported objects are cones of projective characters;
the moves need only transfer numerical bounds, without providing
Morita equivalences.

Throughout this section \(p\) is odd and does not divide \(q\).
We use the ordinary labels and
degree formulas of Proposition~\ref{prop:label-degrees}, the branching
rule of Proposition~\ref{prop:single-cycle}, and the block projections
of Proposition~\ref{prop:core-blocks}.  The type \(A\) series considered
here is unipotent.  In the regular classical groups of
Lemma~\ref{fam:regular-datum}, a \emph{packet} means a Frobenius orbit
of classical factors of the connected centralizer of the fixed semisimple
\(p'\)-parameter whose normalized spectrum is contained in
\(\{1,-1\}\).  We calculate on one representative factor with the
composite endomorphism around the orbit.  There are boundedly many
packets, with field
parameters \(q_i=q^\delta\), \(\delta\leq2\), as in
Lemma~\ref{fam:packets}.  The parameter is retained on every Levi:
induction always starts in its prescribed rational Levi conjugacy
class.

\begin{proposition}\label{prop:runner-reduction}
Fix \(p\) and a defect-order bound \(M\).  Consider either a unipotent
block of \(\GL_n(q)\) or \(\GU_n(q)\) of defect order at most \(M\), or a
block in a sign-parameter series of an original regular classical
group above, before extraction, whose defect order modulo its
central \(p\)-subgroup is at most \(M\).
There are constants depending only on \(p,M\) with the following
properties.

The required Cartan bound reduces to Cartan bounds for groups with
root data in a finite set, and for the following block-union
projections: labels in a fixed ordinary cuspidal triangle series,
with a bounded number of split Harish--Chandra rank-one steps above
the triangle.  The latter projections have boundedly many basic
rows and bounded block defects, although the triangle rank can be
unbounded.  The transfer constants are uniform.

Consequently Theorem~\ref{thm:geometric-cartan} and
Proposition~\ref{prop:triangle-cartan} give the required Cartan bounds
for all these blocks.
\end{proposition}

We first bound the runner weights and calculate the move matrices,
including their type-D foldings.  We then control the complete
reduction in two stages: a map of bounded condition number on the
full coordinate space, followed by the symmetric stage.  Bounds
return through these stages in reverse order, using the two forms
of the cone lemma below.
All dimensions and norms
below are ordinary Euclidean ones; changing to a coordinate
\(\ell^1\)-norm changes constants only by the bounded dimension.

\subsection{Charged runners and bounded weight}

A charged bead set is a subset \(B\subset\Z\) containing every
sufficiently negative integer and no sufficiently positive integer.
Its charge is
\[
 c(B)=|B\cap\Z_{\geq0}|-|\Z_{<0}\setminus B|.
\]
List its elements in decreasing order as \(b_1>b_2>\cdots\).
There is a unique partition \(\lambda\) such that
\[
 b_a=c(B)-a+\lambda_a.
\]
Its weight \(w(B)=|\lambda|\) is also the number of pairs
\((b,g)\) with \(b\in B\), \(g\notin B\), and \(g<b\).
Indeed the bead \(b_a\) has exactly \(\lambda_a\) gaps below it.
The packed set of charge \(c\) is
\(B(c)=\{k\in\Z:k<c\}\).

\begin{lemma}\label{run:localization}
If \(w(B)\leq W\), then \(B\) agrees with \(B(c(B))\) outside
\([c(B)-W,c(B)+W)\).  For charged sets \(B,B'\) of total weight at
most \(W\), put \(m=c(B)-c(B')\).  If \(m>2W\), then
\(B'\subset B\) and \(|B\setminus B'|=m\).  In general
\begin{equation}\label{run:availability}
 |B\setminus B'|\leq \max(m,0)+2W.
\end{equation}
There are at most as many charged sets of fixed charge and weight
\(w\) as partitions of \(w\).
\end{lemma}

\begin{proof}
For \(a>W\) we have \(\lambda_a=0\), and
\(b_1\leq c(B)-1+W\).  This proves the asserted strip.
One can obtain \(B\) from its packed set by \(w(B)\) unit bead
moves; consequently each of the numbers of added beads and removed
beads is at most \(w(B)\).
When \(m>2W\), every gap of \(B\) is above every possible bead
of \(B'\) not already in the common filled negative part, giving
the inclusion.  The difference of charges is then the size of the
set difference.  Comparing both sets with their packed sets proves
\eqref{run:availability}.  The partition description proves the
last assertion.
\end{proof}

To pass from the finite labels to charged bead sets, use
\(\{\lambda_a-a:a\geq1\}\) for a partition \(\lambda\);
this set has charge zero.  For a symbol \((A,B)\), first extend
its rows to \(A\cup\Z_{<0}\) and \(B\cup\Z_{<0}\), of charges
\(|A|\) and \(|B|\).  Translate both sets by
\(-\lfloor(|A|+|B|)/2\rfloor\).  Their charge sum is then zero
or one.  A simultaneous symbol shift translates both extended
sets by one and increases this floor by one, so the normalized
charged rows are independent of the finite representative.

Here are the runner conventions, including the offsets at their
ends.  Each runner records the integers \(k\) for which the
displayed position is occupied in these charged sets.
Write \(\epsilon=1\) in the linear case and \(\epsilon=-1\)
in the unitary case, and set \(e=\ord_p(\epsilon q)\).
Partition positions \(j+ek\) form runner \(j\), with bead set \(B_j\)
in the coordinate \(k\) and packed threshold \(C_j=c(B_j)\).
Extend the indices by
\begin{equation}\label{run:a-offset}
 B_{j+e}=B_j-1,\qquad C_{j+e}=C_j-1.
\end{equation}
A split rank-one extraction moves \(j\) to \(j-1\) for
\(\GL\), and \(j\) to \(j-2\) for \(\GU\), at unchanged \(k\).
Thus the unitary move is a \(2\)-hook removal.

For a symbol, put \(Q=q_i\).
If \(\ord_p(Q)=d\) is odd, use in each row the runners
\(j+dk\), with \(C_{j+d}=C_j-1\).  This is the \emph{hook case}.
If \(\ord_p(Q)=2d\), use positions \(j+dk\) in the first row
when \(k\) is even and in the second row when \(k\) is odd.
Taking \(0\leq j<2d\) lists every position in the two rows once.
Now
\begin{equation}\label{run:cohook-offset}
 \begin{aligned}
 B_{j+2d}=B_j-2,\qquad C_{j+2d}=C_j-2,\qquad
 \\
 (\tau B)_j=B_{j+d}+1,\quad
 (\tau C)_j=C_{j+d}+1,
 \end{aligned}
\end{equation}
where \(\tau\) exchanges the two symbol rows.
This is the \emph{cohook case}.  In both symbol cases a split
rank-one extraction is \(j\longrightarrow j-1\).
All translations in these formulas translate the whole bead set.

Thus partitions have total charge zero, and the two physical symbol
rows have charge sum zero or one.  In types \(B,C\) orient the
rows so that their length difference is \(1\) modulo \(4\).
Packing cohook runners can change this difference by \(2\) modulo
\(4\) at each removal.  The oriented packed core is therefore
determined by the core class and the weight.  Type \(D\) cores are
instead taken up to \(\tau\).  A core is called \emph{symmetric}
if it is fixed by \(\tau\), with the offsets in
\eqref{run:cohook-offset} included.

The sum \(W=\sum_jw(B_j)\), over all the relevant runners and
packets, is the cycle weight.  Packing gives exactly the cores in
Proposition~\ref{prop:core-blocks}.  A hook or cohook at a contributing
distance is an inversion on one of these runners.

\begin{lemma}\label{run:defect-bounds}
For the starting blocks in Proposition~\ref{prop:runner-reduction},
the cycle weights are bounded in terms of \(p,M\).
If all weights vanish, the block is of defect zero after the
indicated central quotient.  Otherwise the relevant basic
cyclotomic \(p\)-parts, and in the regular classical case the actual
central \(p\)-order, are bounded.

In the positive-total-weight case, every endpoint core collection
retaining these weights has boundedly many basic rows and bounded
block defects upstairs.
These bounds also hold for a union of a bounded number of such
collections.
\end{lemma}

\begin{proof}
Every contributing inversion supplies a factor divisible by \(p\)
in the degree-defect formula of
Proposition~\ref{prop:label-degrees}.  In the regular classical
case divide first by the original central \(p\)-order \(|Z_p|\).  Hence
\(p^W\leq M\).
If \(W=0\), every basic-row degree has relative defect one;
the integral-basic-set degree criterion of
Lemma~\ref{lem:row-projection} gives defect zero in the quotient.
This case is completed before any runner extraction.  The toral
\(D_1\) core boundary contributes no hidden \(p\)-factor, by
the explicit packed-core check in Section~\ref{sec:labels}.

If a packet has positive weight, its degree defect is divisible by
the basic factor
\[
 \alpha=
 \begin{cases}
 ((\epsilon q)^e-1)_p,&\text{type }A,\\
 (Q^d-1)_p,&\text{hook case},\\
 (Q^d+1)_p,&\text{cohook case}.
 \end{cases}
\]
Thus the relevant \(\alpha\) is bounded.  All classical packets
have the same \(Q\).  Lemma~\ref{fam:central-order} then bounds
the central \(p\)-order, so that from now on we can work upstairs.
Subsequent residual degree calculations use the inherited torus
and retain any \(D_1\) packet separately, as in
Proposition~\ref{prop:label-degrees}.
Lemma~\ref{fam:central-order} also bounds the full residual
centers, with exponent at most four.  All further steps are
performed upstairs; they do not require a product decomposition
of a quotient by a diagonally acting central subgroup.
Packets of weight zero contribute no such factor.

If an inversion has runner distance \(h\), it accounts for at
least \(h\) inversions: between its gap \(g\) and bead \(b=g+h\),
each intermediate position is either a bead paired with \(g\)
or a gap paired with \(b\).  Therefore \(h\leq W\).
Lifting the exponent, for odd \(p\), bounds its factor by
\(\alpha h_p\); the same statement for a negative cycle follows
by applying it to odd powers in \(Q^d+1\), with the even powers
recorded as hooks on the same alternating runner.
There are at most \(W\) factors.  Thus every basic-row degree
defect is bounded at a retained-weight endpoint.
The integral-basic-set criterion bounds the block defects of the
whole collection, not just the defects of its displayed rows.

There are boundedly many runners, because \(e,d\leq p-1\), and
boundedly many packets.  Lemma~\ref{run:localization} bounds the
number of labels of total weight \(W\), including the at most two
split signs on a type \(D\) label.  The number of blocks is no
greater than the number of basic rows.  The block-union version
of Lemma~\ref{lem:row-projection} now gives all the stated bounds.
\end{proof}

The zero-total-weight case is now complete, so henceforth \(W>0\).
For a retained-weight core collection \(\mathcal C\), let
\(D_{\mathcal C}\) be the matrix whose columns are its PIM
characters projected onto its ordinary basic rows.  The lemma
makes this a square, full-rank nonnegative integral matrix of
bounded size and bounded determinant.  Its column cone is
\[
 \mathcal P_{\mathcal C}
   =D_{\mathcal C}\mathbb R_{\geq0}^{\,l(\mathcal C)},
\]
where \(l(\mathcal C)\) is the number of simple modules in the
block collection.  We will use induction to compare these cones
and thereby bound the entries of \(D_{\mathcal C}\).
Lemma~\ref{lem:row-projection} then bounds the full projective
characters and their Cartan inner products.

\subsection{Ordered moves and their realization}

An edge means one of the allowed moves \(j\to j'\).
Its positive difference is \(C_j-C_{j'}\), using the displayed
offsets if an index passes an end.  If this difference is
\(m>2W\), the corresponding bead sets are nested.  Transfer the
\(m\) beads in \(B_j\setminus B_{j'}\) across the edge.
This interchanges the two bead sets, with their coordinate offsets.

\begin{lemma}\label{run:ordered-swap}
Suppose the ordered core convention is allowed at both ends.
The projection of \(m\) successive rank-one extractions onto the
swapped core has matrix \(m!P\), where \(P\) is the bijection
interchanging the runner bead sets.  The reversed induction matrix
is its transpose.  The same conclusion holds for two disjoint
edges, each with difference \(m\), with coefficient \((2m)!\)
and the simultaneous-swap bijection.
\end{lemma}

\begin{proof}
Let \(v_j\) count the moves across the edge from \(j\).
Their effect on the charges is the incidence vector \(\partial v\);
on a step-one cycle,
\[
 (\partial v)_j=v_{j+1}-v_j.
\]
Changes of charges have no offset.  The kernel of this incidence
map consists precisely of vectors constant on each directed cycle.
The prescribed charge change is \(\partial(me_j)\).
The difference between any other flow and \(me_j\) is therefore
constant on each cycle.  Each cycle has an unused edge, so
nonnegativity forces every such constant to be nonnegative.
The total number \(m\) of extractions forces every constant to be
zero.  For two disjoint opposite edges use
\(me_j+me_{j^*}\) instead.  In the cases where we use this
argument each cycle again has an unused edge.

Thus a bead can move only across the prescribed edge or edges.
It can cross each edge at most once.  Every surplus bead must
cross, and any order of the \(m\), respectively \(2m\), distinct
transfers is legal.  This gives the factorial and the unique
output.  The total runner weight is unchanged by permuting the
bead sets.  Conversely every target label has the same bounded
weight, so Lemma~\ref{run:localization} supplies the reverse
nesting and reconstructs its unique source.  Path reversal proves
the assertion for induction.
\end{proof}

The cycles omitted by the unused-edge assertion cannot have a
large positive difference: for a length-one cycle the difference
is the negative period offset.  In particular this covers
\(e=1\) in \(\GL\), \(e\in\{1,2\}\) in \(\GU\), and \(d=1\)
in hook symbols.

Each individual transfer is an actual hook removal on a partition
or a same-row hook removal on a symbol, and therefore an actual
split Harish--Chandra step.  By Lemma~\ref{fam:packets}, a
classical packet step extracts a factor \(\GL_\delta(q)\);
the type \(A\) extracts are \(\GL_1(q)\) and \(\GL_1(q^2)\).
All these factors have order prime to \(p\) whenever a large
move occurs.  In the hook case a large move requires \(d>1\);
in the cohook case \(p\nmid Q-1\).  If \(\delta=2\), the former
condition also excludes \(p\mid q^2-1\), and the latter does so
directly.  The type \(A\) exclusions above likewise give the claim.
These extracted factors can be retained as independent frozen
factors in the Levi, each in its fixed defect-zero block.

There is no additional rank or sign constraint to impose after a
legal path has been specified.  Ordinary Harish--Chandra branching
realizes each of its labels in the appropriate residual classical
group.  The physical row lengths, and thus the symbol defect
congruence, are preserved by each symbol step.  A nonsplit
orthogonal residual group cannot reach a rank-zero label.
The independent hyperbolic coordinates in
Lemma~\ref{fam:packets} realize the whole chain of extractions,
including the fixed other packets.

Induction and restriction are exact and preserve projectives:
the unipotent radical has order prime to \(p\), so its averaging
idempotent defines an exact direct summand.  Projection onto a
union of blocks also preserves projectives.  All resulting
projective maps are therefore nonnegative in PIM bases.
By the ordinary parameter rule in
Proposition~\ref{prop:single-cycle}, rows with a nontrivial
semisimple \(p\)-part cannot contribute to the basic rows with
trivial \(p\)-part.  Consequently their projected coordinate maps
depend only on the displayed basic coordinates.  This remains
true with intermediate core projections.  Possible other Levi
parameter classes inducing to the same ambient class are
irrelevant, since the source has been projected onto the
prescribed class.

For an unnormalized composite functor \(F\) between retained-weight
core collections \(\mathcal C\) and \(\mathcal C'\), let
\(A_F\) be this map on basic coordinates.  Let \(N_F\) record
the multiplicities of target PIMs in the images of source PIMs.
The two descriptions of each image give
\[
 A_FD_{\mathcal C}=D_{\mathcal C'}N_F,
 \qquad N_F\geq0,
 \qquad
 A_F\mathcal P_{\mathcal C}\subseteq\mathcal P_{\mathcal C'}.
\]
The entries of \(N_F\) are integers, but no bound on them is
assumed.  Dividing \(A_F\) by a positive scalar preserves the
cone inclusion, with the same division applied to \(N_F\).
The bounded sizes and determinants concern the retained-weight
endpoints; no bound on intermediate projection dimensions is needed.

\subsection{The type D folding}

The ordered swap calculation is now available as an actual
projective-preserving operation.  In type $D$ its ordinary matrix
must still be checked after row exchange, because an unordered
target can receive extra paths and an equal-row label has two
distinct characters.

\begin{lemma}\label{run:folding}
In a type \(D\) ordinary Harish--Chandra series with equal row
lengths, folding ordered bipartitions sends a nonsymmetric
bipartition to its unordered character and an equal bipartition
to the sum of its two split characters.  With twice the single
character as the folding convention at relative rank zero, this
map intertwines induction through any positive number of marked
rank-one extractions and every underlying-core projection.
Such induction annihilates the difference of two split source
characters.
\end{lemma}

\begin{proof}
Write \(B_b=C_2\wr S_b\) and let \(D_b\) be its subgroup of even
signed permutations.  The asserted folding at positive rank is
\(\Res^{B_b}_{D_b}\), by the elementary index-two character
description of these groups.  If \(1\leq j<b\), embed \(B_j\)
on the old indices and fix every new index.  An odd sign change
on an old index belongs to \(B_j\), so
\[
 D_bB_j=B_b,\qquad D_b\cap B_j=D_j.
 \]
The Mackey formula, with its single double coset, gives
\[
 \Res^{B_b}_{D_b}\Ind^{B_b}_{B_j}
   =\Ind^{D_b}_{D_j}\Res^{B_j}_{D_j}.
 \]
At \(j=0\) there are two double cosets instead, and the left
side is twice \(\Ind^{D_b}_{D_0}\); this proves the stated
rank-zero convention.

An odd sign change on an old index interchanges the two split
characters.  Multiplying it by a sign change on a new index
gives an element of \(D_b\) inducing the same conjugation on
\(D_j\).  Their induced characters are thus equal, proving
the assertion about their difference.
Core projectors retain precisely a row-exchange orbit and both
split signs if present.  They commute with folding.  Iteration
proves the assertion with intermediate projectors.
\end{proof}

For nonzero row-length difference orient by the length and use
ordinary ordered branching.  For zero difference
Lemma~\ref{run:folding} shows that, between nonsymmetric cores,
the coefficient is the count of ordered paths from one fixed
orientation to either orientation of the target.  There are no
equal-row labels above a nonsymmetric core.  Reversing paths and
exchanging both orientations gives the identical rule for
induction.

\begin{lemma}\label{run:nonsymmetric}
Suppose \(d>1\), and let \(j^*\) be the edge paired with \(j\)
by row exchange.  They are disjoint.  A sufficiently large swap
from a nonsymmetric core to a nonsymmetric core has exactly the
matrix of Lemma~\ref{run:ordered-swap}, also on unordered labels.
If \(d\geq3\) and a sufficiently large swap would give a symmetric
core, there is a different large edge whose swap stays
nonsymmetric.
\end{lemma}

\begin{proof}
Let \(m=C_j-C_{j'}\), and put
\(m_1=C_{j^*}-C_{(j')^*}\).
If \(v\) is a flow to the exchanged target, applying row
exchange to its charge equation and adding gives
\[
 \partial(v+\tau v)=\partial(me_j+me_{j^*}).
 \]
Every relevant cycle has a missing edge because \(d>1\)
(in the hook case \(d\) is odd, hence at least three).
Nonnegativity and the total length \(2m\) therefore give
\(v+\tau v=me_j+me_{j^*}\).
In particular \(v\) is supported on these two disjoint edges.
Their charge equations give
\begin{equation}\label{run:exchanged-flow}
 v_j=\frac{m-m_1}{2},\qquad
 v_{j^*}=\frac{m+m_1}{2}.
\end{equation}
If \(|m_1|\leq2W\), the second expression is larger than the
availability bound \eqref{run:availability} for sufficiently
large \(m\).  If \(m_1>2W\), nesting gives
\(v_{j^*}\leq m_1\), whereas \(v_j\geq0\) gives \(m_1\leq m\).
Equality \(m_1=m\) is forced, so \(v_j=0\); the full charge
equation then says that the source is symmetric.  If
\(m_1<-2W\), no move at \(j^*\) is available, so
\(v_{j^*}=0\); the target is symmetric.  Both possibilities
are excluded.

Now suppose that the intended target is symmetric.  The source
and its exchange then differ exactly on the two pairs of
vertices incident to the edge and its opposite.  In the hook
case, put the thresholds on one edge equal to \(L,H\), with
\(m=H-L\).  The other row has \(H,L\) at these vertices.
Along the remaining \(d-1\) edges of that other row, the
threshold rises sum to \(m-1\), the subtraction of one coming
from the period offset.  One of these other edges consequently
has positive difference at least \((m-1)/(d-1)\).

In the cohook case number the vertices so that the four
exceptional thresholds are
\[
 C_0=L,\quad C_1=H,\quad C_d=H-1,\quad C_{d+1}=L-1.
 \]
At every other pair, \(C_{i+d}=C_i-1\).
Put \(\Delta_i=C_i-C_{i-1}\).
For \(d\geq3\), the identity
\[
 \Delta_{d+2}+\sum_{i=3}^{d}\Delta_i=m-1
 \]
again supplies a different positive difference at least
\((m-1)/(d-1)\).
Choose the original threshold large enough that this exceeds
the inclusion threshold.  The new edge is incident to a
different pair of row-exchange vertex pairs.  Its charge
change cannot cancel the original asymmetry, which is supported
on exactly the original two pairs.  Its target is therefore
nonsymmetric.
\end{proof}

\begin{lemma}\label{run:symmetric}
For a symmetric core and \(d>1\), simultaneous swaps on two
opposite large edges have coefficient \((2m)!\) on ordered
labels.  The output bijection commutes with row exchange and
preserves equal-row labels.
After factorial normalization, the composite of any sequence
of these inductions has bounded rational entries and bounded
denominators.  Its image is the subspace in which coordinates
on each affected split pair agree, with the full coordinates
retained on unaffected packets.
\end{lemma}

\begin{proof}
Symmetry makes the two edge differences equal.  The flow proof
of Lemma~\ref{run:ordered-swap} gives the unique simultaneous
flow and the \((2m)!\) possible orders.  Both the operation and
its inverse commute with row exchange, so they preserve fixed
points as well as two-element orbits.

Use a basis consisting of the nonsplit labels and the sums of
split pairs.  Lemma~\ref{run:folding} says that normalized
induction is the resulting bijection on this basis, except
for its factor of two at rank zero.  For an individual split
source its image is half the image of the sum, because the
difference is annihilated as soon as a new index is added.
At each subsequent step the sum, rather than either individual
label, is transported bijectively.  Thus the factor \(1/2\)
occurs only once.  Rank zero is passed at most once in an
induction chain.  These observations also prove that the image
is the entire indicated subspace.  Tensoring over the bounded
number of packets preserves the bounded-entry and
bounded-denominator assertions.  Packets on which no such
induction occurs retain their full space.
\end{proof}

\subsection{The two exceptional cohook moves}

Only cohook \(d=1,2\) remain.  Their normalized matrices need not
be permutations, but the errors form convergent series.  For
\(d=1\), repeated moves of the same bead give an exponentially
small extra contribution.  For \(d=2\), a two-edge path avoids
the symmetric core and gives a quadratic error bound.  Both
estimates will control products of arbitrarily many moves.

\begin{lemma}\label{run:d-one}
For cohook \(d=1\) and sufficiently large positive difference
\(m\), both endpoints of the swap are nonsymmetric.
After division by \(m!\), its unordered matrix is \(P+E_m\),
where \(P\) is the ordered-swap permutation and, for constants
depending only on the weight,
\[
 \|E_m\|\leq K\,2^{-m/2}.
 \]
The next positive difference is \(m-2\).
\end{lemma}

\begin{proof}
Write \(C_0=L,C_1=H,C_2=L-2\), with \(m=H-L\).
The ordered target is \((H,L)\), whereas its exchange is
\((L+1,H-1)\).  A path of length \(m\) to the latter has
exactly \((m+1)/2\) moves on \(1\to0\) and \((m-1)/2\)
on the other edge, by its charge equations.  If these are not
integers, there is no such path.

Below \(L-W-2\) both runners are full.  No bead in this region
can ever move, since a move goes downwards and would require
a gap below it.  There are at most \(K_0=K_0(W)\) beads
originally on the low runner above this region.  Every move
on the other edge is therefore a repeated move of a bead
already moved earlier, except possibly for these \(K_0\)
beads.  At least \((m-1)/2-K_0\) moves repeat a bead.

Fix the two oriented endpoint labels.  Within each physical
row beads cannot pass each other: all these moves have
physical length one.  Matching them in their order fixes
each bead's number \(a_t\) of moves.  The number of possible
paths is at most
\[
 \frac{m!}{\prod_t a_t!}.
 \]
For \(a\geq1\), \(a!\geq2^{a-1}\).  Hence
\(\prod_ta_t!\geq2^{(m-1)/2-K_0}\).
This bounds every extra normalized entry exponentially.
The endpoint dimensions are bounded by
Lemma~\ref{run:defect-bounds}, giving the claimed operator
bound.

The source is symmetric only for \(m=-1\), and the swapped
target only for \(m=1\), so neither large endpoint is symmetric.
After the swap the other edge has
positive difference \(H-(L+2)=m-2\), as claimed.
\end{proof}

\begin{lemma}\label{run:d-two}
In cohook \(d=2\), suppose a large swap from a nonsymmetric
core would give a symmetric core.  Number its thresholds as
\[
 [C_0,C_1,C_2,C_3]=[L,H,H-1,L-1],\qquad m=H-L.
 \]
There is a projected extraction of length \(2m-1\) to the
nonsymmetric core
\([H,H-1,L,L-1]\).
After division by \(m!(m-1)!\), its matrix has the form
\[
 (1-1/m)P+E_m,\qquad
 \|E_m\|\leq \frac{K(W)}{m(m-1)},
 \]
where \(P\) is a bijection preserving total runner weight.
Reindexing the target from the old vertex \(3\) gives the
same exceptional pattern with parameter \(m-1\).
\end{lemma}

\begin{proof}
Use the ordered edge sequence
\begin{equation}\label{run:detour}
 (1\to0)^{m-1},\quad (2\to1),\quad
 (1\to0),\quad (2\to1)^{m-2},
\end{equation}
projecting onto its core after every individual step.
After \(u,v\) steps on the two edges the thresholds are
\[
 [L+u,H-u+v,H-1-v,L-1].
 \]
The symmetry equations are \(u+v=m\) and \(u-v=m\);
thus symmetry occurs exactly at \((u,v)=(m,0)\).
The path \eqref{run:detour} avoids it.

A single step between these nonsymmetric cores cannot reach
the opposite orientation.  Indeed adding its flow to its
row exchange leaves total length two, supported on that edge
and its opposite by the missing-edge argument.
Choosing the same edge for the exchanged target would make
the target symmetric; choosing the opposite edge would make
the source symmetric.  Both are excluded.  Thus the actual
unordered projections force exactly the ordered trajectory,
even when a particular label has equal row lengths.

The initial nesting gives \(B_0\subset B_1,B_2\).
Set
\[
 A=B_1\setminus B_0,\qquad C=B_2\setminus B_0,\qquad
 |A|=m,\quad |C|=m-1,\quad \rho=|C\setminus A|.
 \]
Here \(\rho\leq W\): every bead of \(B_2\setminus B_1\)
is either a hole below the threshold of \(B_1\) or an extra
bead above the threshold of \(B_2\), and these holes and
extra beads are bounded by the two runner weights.

The first segment transfers all of \(A\) except one choice
\(a\).  The next step transfers a choice
\(b\in C\setminus\{a\}\).
If the following step transfers \(a\), the last segment must
transfer all remaining elements of \(C\).  Its output is
the rotation
\[
 [B_0,B_1,B_2]\longmapsto[B_1,B_2,B_0].
 \]
For each such pair \((a,b)\), the first and last long
segments can be ordered in
\((m-1)!(m-2)!\) ways.  The number of pairs is
\(m(m-1)-|A\cap C|\).  The exact rotation coefficient is
therefore
\begin{equation}\label{run:rotation-count}
 m!(m-1)!
 \left(1-\frac{|A\cap C|}{m(m-1)}\right).
\end{equation}

The only alternative is to transfer \(b\) at the third
segment.  It can enter runner \(0\) only if \(b\notin A\).
Completing all \(m-2\) final transfers is then possible only
if \(a\notin C\); otherwise one fewer bead can enter runner
\(1\).  Conversely these two conditions suffice, and give
exactly the same number \((m-1)!(m-2)!\) of paths.
There are \(\rho(\rho+1)\) such pairs, since
\(|A\setminus C|=\rho+1\).
Also \(|A\cap C|=m-1-\rho\).  Thus the column \(\ell^1\)-error
from \((1-1/m)P\), after normalization, is exactly
\[
 \frac{\rho+\rho(\rho+1)}{m(m-1)}
   =\frac{\rho(\rho+2)}{m(m-1)}
   \leq\frac{W(W+2)}{m(m-1)}.
 \]
The bounded endpoint dimension gives the operator estimate.

The rotation is a bijection: at the target its low runner
is still nested in both high runners, so the inverse rotation
recovers every endpoint label.  It permutes runner
partitions and hence preserves their total weight.  Each
transfer uses the same physical row at its source and
target; physical row lengths are preserved.
Finally the period offset gives, on reindexing from \(3\),
\[
 [L-1,H-2,H-3,L-2],
 \]
which has parameter \(m-1\).  No bound on the dimensions of
the intermediate projections was used.
\end{proof}

\subsection{Termination and transfer of bounds}

All individual move matrices have now been calculated.  We must
choose a terminating sequence and control its cumulative effect:
bounded condition number handles the invertible stage, while a
separate argument handles the split-pair identifications.

Choose all large-difference thresholds in terms of \(p,W\),
larger than the inclusion thresholds above.  At a nonsymmetric
type \(D\) core use an exact swap when its target is nonsymmetric.
For \(d\geq3\), Lemma~\ref{run:nonsymmetric} replaces a forbidden
swap by a different large swap.  Stop once every positive
difference is below the resulting fixed threshold.
For cohook \(d=1\) use Lemma~\ref{run:d-one}; for cohook \(d=2\),
use exact swaps until stopping or first reaching the exceptional
pattern, then use Lemma~\ref{run:d-two} repeatedly.
All other ordered packets use Lemma~\ref{run:ordered-swap}.
Perform these full-space reductions before reducing symmetric
packets by the paired moves of Lemma~\ref{run:symmetric}.

Every move removes positive actual rank, so the process terminates.
This also proves termination when replacing a forbidden edge
by a different one.  The special cohook sequences have strictly
decreasing parameters \(m\), by two or by one, respectively.
There is at most one such sequence of each kind per packet:
once the exceptional pattern is reached its own rule continues
to the stopping threshold.  Exact moves between any other
stages have permutation matrices after normalization.

Number the basic-coordinate spaces at the starting collection,
the end of the full-space stage, and the final endpoint by
\(0,1,2\), and write \(n_j\) for their respective dimensions.
Reversing the extractions gives the normalized induction maps
\[
 \mathbb R^{n_2}\xrightarrow{\ A\ }\mathbb R^{n_1}
              \xrightarrow{\ T\ }\mathbb R^{n_0}.
\]
Bounds will return along these arrows.  The full-space map \(T\)
will be controlled by its condition number.  The map \(A\) can
annihilate split-pair differences.  Its bounded entries and
denominators, together with its image containing the vector of
ordinary character degrees, will instead control the PIM columns
at stage \(1\).

\begin{lemma}\label{run:condition}
The normalized ordinary-coordinate induction from the endpoint
of the full-space stage to its starting point is invertible,
with uniformly bounded condition number.
\end{lemma}

\begin{proof}
Normalize an exact move by its factorial.  Normalize a
cohook-\(1\) move in the same way, and a cohook-\(2\) move
by \(m!(m-1)!(1-1/m)\).
Every matrix is then a permutation plus an error of norm
\(\varepsilon_m\), where the sums of all errors are bounded
uniformly: they are bounded by a constant times
\(\sum_{m\geq m_*}2^{-m/2}\) or
\(\sum_{m\geq m_*}m^{-2}\), where \(m_*\) is the stopping cutoff.
Take \(m_*\) large enough that every \(\varepsilon_m<1/2\).
The singular values of a permutation plus that error lie
between \(1-\varepsilon_m\) and \(1+\varepsilon_m\).
Consequently the condition number of the product is at most
\[
 \prod_m\frac{1+\varepsilon_m}{1-\varepsilon_m}
 \leq \exp\!\left(4\sum_m\varepsilon_m\right).
 \]
The bounded number of packets gives a uniform bound.
\end{proof}

Here is the precise linear algebra that converts these maps to
bounds on projective characters.  It is useful that a scalar
normalization, however small, preserves projective positivity.

\begin{lemma}\label{run:cone-transfer}
Let \(V_-,V_+\) be nonnegative integral square invertible
matrices of bounded size, with
\(|\det V_+|\leq\Delta\), and suppose the entries of \(V_-\)
are bounded.

If an invertible map \(T\) of bounded condition number sends
the cone generated by the columns of \(V_-\) into the cone
generated by those of \(V_+\), the entries of \(V_+\) are bounded.

Alternatively let \(A\), possibly rectangular, have bounded
rational entries with a common bounded denominator.  Suppose
\[
 V_+^{-1} A V_-\geq0
 \]
entrywise, and the image of \(A\) contains a vector
\(V_+d\) with every coordinate of \(d\) strictly positive.
Then the entries of \(V_+\) are bounded.
\end{lemma}

\begin{proof}
Write \(\mathcal C_\pm\) for the two cones.
There are only finitely many possible bounded integral
matrices \(V_-\) in each of the bounded dimensions.
Thus the volume of \(\mathcal C_-\) in the unit ball has a
positive uniform lower bound.  Rescale \(T\) so its largest
singular value is one.  If its condition number is at most
\(\kappa\), its determinant in absolute value is at least
\(\kappa^{-n}\).  Its image of this part of
\(\mathcal C_-\) lies in \(\mathcal C_+\) and in the unit
ball.  The corresponding volume for \(\mathcal C_+\)
therefore also has a positive uniform lower bound.
Since all coordinates in the latter cone are nonnegative,
its cut by \(\sum x_i\leq1\) has a positive uniform
volume lower bound as well.

If \(s_j\) is the coordinate sum of column \(j\) of \(V_+\),
the linear change of variables \(x=V_+t\) gives exactly
\[
 \operatorname{vol}
 \{x\in\mathcal C_+:\textstyle\sum_i x_i\leq1\}
   =\frac{|\det V_+|}{n!\prod_j s_j}.
 \]
Every \(s_j\) is a positive integer.  The upper bound
\(\Delta\) for the numerator bounds their product and
hence each \(s_j\), proving the first assertion.

For the second, put \(M_0=V_+^{-1}AV_-\).
Because \(V_-\) is invertible, its image is
\(V_+^{-1}\operatorname{im}A\), which contains \(d\).
No row of \(M_0\) can therefore be zero.
Nonnegativity makes every entry of
\(\lambda=M_0\mathbf1\) positive.
If \(Q\) is a common denominator for \(A\), Cramer's rule
shows that the denominator of each entry of \(\lambda\)
divides \(Q\det V_+\).  In particular
\(\lambda_j\geq1/(Q\Delta)\).
The vector \(y=AV_-\mathbf1\) has bounded entries, and
\[
 y=V_+\lambda
 \]
is a nonnegative combination of every column of \(V_+\)
with these uniformly positive coefficients.  It bounds
each column entry, proving the claim.
\end{proof}

Apply the second assertion first, to transfer a bound from the
end of the symmetric stage to its beginning.  Its matrix \(A\)
is the one in Lemma~\ref{run:symmetric}.  Its image contains
the vector of ordinary degrees on the basic rows, since
those degrees agree on split signs by
Proposition~\ref{prop:label-degrees}.  This vector is the
projection of the regular character and equals
\(V_+d\), where \(d\) lists the positive dimensions of the
simple modules.  Thus all hypotheses of the second assertion
hold.  The bounded determinant is supplied by
Lemma~\ref{run:defect-bounds}; no bound on \(V_+^{-1}\)
is assumed.

Next apply the first assertion to the full-space stage, using
Lemma~\ref{run:condition}.  Projective positivity gives the
required inclusion of cones.  Thus a bound at the final
endpoint gives a bound at the original core collection.
Lemma~\ref{lem:row-projection} then gives the full
projective-character and Cartan bounds.  These arguments
require no equivalence of module categories associated to a
runner move.

\subsection{The endpoint and its triangle series}

At the stopping point, every positive edge difference is
bounded.  On each directed cycle the sum of differences is
its fixed period offset, with a minus sign.  A lower bound
for each difference follows by summing the upper bounds
for the others.  Thus all thresholds in a cycle have bounded
spread.

If the runner graph has a single cycle, fixed total charge
then bounds every threshold.  The packed configurations have
bounded rank, and the bounded runner weights give bounded
rank for every endpoint label.  This applies to \(\GL\),
to \(\GU\) with \(e\) odd, and to cohook symbols.
In the regular classical case the packets are all of the
same hook or cohook kind, because their \(q_i\)'s agree.
The accumulated frozen factors split off by
Lemma~\ref{fam:packets}.  The remaining root data belong to
a finite set by Lemma~\ref{fam:regular-datum}: restricting
the regular lattice after extraction retains its uniform
bounded-denominator coordinate description, now in bounded
dimension.  The geometric theorem applies to that finite
set, while every frozen block is of defect zero.

There are two remaining graphs.  If \(e=2d\) in \(\GU\),
the step-two graph has separate parity cycles.  In the hook
symbol case there is one cycle in each physical row.
Their mean thresholds can differ without bound.  This
difference records precisely the ordinary cuspidal triangle:
the \(2\)-core staircase of a unitary partition, or the
symbol obtained by packing its two rows separately at
step one, taken up to row exchange in type \(D\).

\begin{lemma}\label{run:triangle-endpoint}
At these two-cycle endpoints the number \(m_i\) of ordinary
split Harish--Chandra rank-one steps above the triangle in
each packet is uniformly bounded.
The complete collection at this triangle and this \(m_i\)
is a union of the core block projections of
Proposition~\ref{prop:core-blocks}.  It has boundedly many
basic rows and bounded block defects, independently of
the triangle ranks.
\end{lemma}

\begin{proof}
Consider one physical row in the hook case, or one parity
of positions in the unitary case.  Translate all its runner
thresholds by a common integer until one is zero.  This
translation does not affect the number of inversions for
the ordinary step, which is one or two respectively.
Bounded spread and Lemma~\ref{run:localization} put every
deviation from a fully packed row or parity inside an
interval of bounded length.  Outside that interval all
lower positions are occupied and all upper positions
empty.  Every inversion has both ends in the interval,
so their number is bounded by the square of its length.
Summing over the two rows or parities bounds \(m_i\).
Packing at the ordinary step is exactly removal to the
ordinary cuspidal triangle, so this inversion count is
its ordinary Harish--Chandra weight.

A same-row \(d\)-hook preserves the row lengths of a
symbol and hence its triangle.  An \(e\)-hook in the
unitary case has even length and preserves its \(2\)-core.
Thus a whole relevant core class lies in one triangle
series.  At the fixed ambient packet rank, \(m_i\) is
determined by that triangle; enlarging to all labels with
these data merges complete core classes and is therefore
still a projection onto a union of blocks.

The ordinary branching parametrizes these labels by
partitions or bipartitions of the bounded integers \(m_i\),
with the stated split signs in type \(D\).  Their number
is bounded independently of the triangle.
Every contributing relevant inversion is an ordinary-step
inversion, so there are at most \(m_i\), and its distance
in units of the relevant step is at most \(m_i\).
The already bounded basic cyclotomic \(p\)-part and lifting
the exponent bound all their degree defects.  The central
\(p\)-order is bounded as in
Lemma~\ref{run:defect-bounds}.  Applying the integral-basic-set
criterion once more bounds the defects of every block in
this enlarged collection.
\end{proof}

\begin{proof}[Proof of Proposition~\ref{prop:runner-reduction}]
Choose the core collection containing the basic rows of the
starting block.  It is a union of blocks by
Proposition~\ref{prop:core-blocks}.
Lemma~\ref{run:defect-bounds} handles total weight zero directly
and supplies all uniform defect and dimension bounds otherwise.
The terminating runner process and the two forms of
Lemma~\ref{run:cone-transfer} reduce the problem to the endpoints
just described.  A single-cycle endpoint is covered by
Theorem~\ref{thm:geometric-cartan} and the finite-root-data
argument above.

The remaining endpoint in Lemma~\ref{run:triangle-endpoint} is exactly
the endpoint of Proposition~\ref{prop:triangle-cartan}.
In the unitary case its notation is
\(Q_i=q^2\), \(d=\ord_p(Q_i)\), \(e=\ord_p(-q)=2d\);
in the hook classical case it is
\(Q_i=q_i\), \(d=\ord_p(Q_i)\) odd.
The frozen factors can either be retained in their
defect-zero blocks or removed using the Levi product.
The endpoint proposition supplies the remaining uniform
bound.  A Cartan bound bounds every projected decomposition
entry, since a Cartan diagonal entry is the sum of the
squares of the corresponding full column.
The two transfers proved above carry it back
through all runner moves, and
Lemma~\ref{lem:row-projection} supplies the Cartan bound.
For a central quotient we finally use
Lemma~\ref{lem:central-transfer}; the weight-zero case was
already treated directly in the quotient.
This proves Proposition~\ref{prop:runner-reduction}.
\end{proof}
